\chapter{Sprints 6 et 7~: Projets, Reporting et Analyse Prédictive}

\noindent\textbf{Plan du chapitre}
\begin{enumerate}[nosep]
  \item Sprint 6~: Gestion des projets, tâches et scoring de CV
  \item Sprint 7~: JasperReports, tableaux de bord et module d'attrition
\end{enumerate}

% ─────────────────────────────────────────────────────────────
\section{Introduction}
% ─────────────────────────────────────────────────────────────

La troisième release ajoute deux dimensions à Synapse~: une dimension
\textit{collaborative} avec la gestion des projets et des tâches, qui dote les
chefs d'une visibilité en temps réel sur l'avancement sans saisie manuelle
d'indicateurs~; et une dimension \textit{décisionnelle} avec le reporting
JasperReports et le module de prédiction de l'attrition, qui transforme les
données accumulées en signaux exploitables par la direction RH.

% ─────────────────────────────────────────────────────────────
\section{Sprint 6~: Projets, tâches et scoring de CV}
% ─────────────────────────────────────────────────────────────

\subsection{Objectifs}

\begin{table}[H]
\centering
\caption{Périmètre du sprint 6}
\label{tab:s6}
\begin{tabularx}{\textwidth}{|l|L|}
\hline
\rowcolor{headerblue!25}\textbf{Tâche} & \textbf{Résultat attendu} \\\hline
\rowcolor{rowgray}
\textit{projets-service}~(:8087) & CRUD projets, affectation multi-membres, calcul
  automatique du pourcentage d'avancement, dates jalons \\\hline
\textit{taches-service}~(:8088) & CRUD tâches avec statuts À\_FAIRE / EN\_COURS /
  TERMINEE / EN\_RETARD, affectation à un membre, filtre par statut \\\hline
\rowcolor{rowgray}
Kanban Angular & Colonnes dynamiques par statut, drag-and-drop simulé par boutons
  de transition, indicateurs de délai dépassé \\\hline
Recrutement + scoring CV & Formulaire d'offre, dépôt de CV, score automatique
  calculé par Groq sur 100 avec justification structurée \\\hline
\end{tabularx}
\end{table}

\subsection{Diagramme de cas d'utilisation~: projets et tâches}

\begin{figure}[H]
\centering
\begin{tikzpicture}

  % Acteur Chef
  \draw[thick] (-2,5.5) circle(0.28cm);
  \draw[thick] (-2,5.22) -- (-2,4.5) coordinate(wC);
  \draw[thick] (wC)++(0,0.1)--++(-0.35,-0.3) (wC)++(0,0.1)--++(0.35,-0.3);
  \draw[thick] (wC)--++(-0.25,-0.5) (wC)--++(0.25,-0.5);
  \node[font=\small] at (-2,3.7) {Chef};

  % Acteur Employé
  \draw[thick] (12,3.5) circle(0.28cm);
  \draw[thick] (12,3.22) -- (12,2.5) coordinate(wE);
  \draw[thick] (wE)++(0,0.1)--++(-0.35,-0.3) (wE)++(0,0.1)--++(0.35,-0.3);
  \draw[thick] (wE)--++(-0.25,-0.5) (wE)--++(0.25,-0.5);
  \node[font=\small] at (12,1.7) {Employé};

  % Frontière
  \begin{scope}[on background layer]
    \filldraw[umlborder] (0,0) rectangle (10,7.5);
    \node[font=\bfseries\small] at (5,7.2) {projets-service / taches-service};
  \end{scope}

  % UCs
  \draw (4,6.5) ellipse(2.4cm and 0.45cm); \node[font=\scriptsize,align=center] at (4,6.5) {Créer un projet};
  \draw (7.5,6.5) ellipse(2cm and 0.45cm); \node[font=\scriptsize,align=center] at (7.5,6.5) {Affecter des membres};
  \draw (4,5.1) ellipse(2.4cm and 0.45cm); \node[font=\scriptsize,align=center] at (4,5.1) {Suivre l'avancement (\%)};
  \draw (4,3.7) ellipse(2.4cm and 0.45cm); \node[font=\scriptsize,align=center] at (4,3.7) {Créer une tâche};
  \draw (7.5,3.7) ellipse(2cm and 0.45cm); \node[font=\scriptsize,align=center] at (7.5,3.7) {Assigner une tâche};
  \draw (4,2.3) ellipse(2.4cm and 0.45cm); \node[font=\scriptsize,align=center] at (4,2.3) {Mettre à jour le statut};
  \draw (4,0.9) ellipse(2.4cm and 0.45cm); \node[font=\scriptsize,align=center] at (4,0.9) {Consulter les retards};

  % Relations
  \draw[->,dashed,>=stealth] (6.4,6.5) -- (5.5,6.5) node[midway,above,font=\scriptsize]{<<include>>};
  \draw[->,dashed,>=stealth] (5.5,3.7) -- (5.4,3.7) node[midway,above,font=\scriptsize]{<<include>>};

  % Associations acteurs
  \draw[->] (-1.65,5.5) -- (1.6,6.5);
  \draw[->] (-1.65,5.3) -- (1.6,5.1);
  \draw[->] (-1.65,5.0) -- (1.6,3.7);
  \draw[->] (-1.65,4.7) -- (1.6,0.9);
  \draw[->] (11.65,3.5) -- (6.4,2.3);
  \draw[->] (11.65,3.2) -- (6.4,0.9);

\end{tikzpicture}
\caption{Diagramme de cas d'utilisation~: gestion des projets et tâches}
\label{fig:uc_proj}
\end{figure}

\subsection{Réalisation~: calcul d'avancement}

L'avancement d'un projet est recalculé dans \textit{projets-service} à chaque
mise à jour de statut d'une tâche~:

\begin{lstlisting}[language=Java,caption={Calcul de l'avancement dans ProjetService}]
double avancement = taskRepo
    .countByProjetIdAndStatut(projetId, "TERMINEE")
    * 100.0
    / taskRepo.countByProjetId(projetId);
projet.setAvancement(Math.round(avancement));
projetRepo.save(projet);
\end{lstlisting}

\subsection{Réalisation~: scoring de CV par Groq}

Pour chaque candidature déposée, le \textit{chat-service} construit un prompt
intégrant le descriptif du poste et le texte extrait du CV~; il appelle
\texttt{llama-3.3-70b-versatile} avec \texttt{temperature=0.1} (résultat
déterministe) et demande un objet JSON \texttt{\{score: int, forces: [...],
faiblesses: [...], recommandation: string\}}. Le score entier (0--100) et la
recommandation sont persistés dans la table \texttt{CANDIDATURES}.

\subsection{Résultats du sprint 6}

\begin{figure}[H]
  \centering
  % \includegraphics[width=\textwidth]{images/sprint6_kanban.png}
  \caption{Tableau Kanban d'un projet~: colonnes À\_FAIRE / EN\_COURS / TERMINÉE / EN\_RETARD}
  \label{fig:s6_kanban}
\end{figure}

% ─────────────────────────────────────────────────────────────
\section{Sprint 7~: JasperReports, tableaux de bord et attrition}
% ─────────────────────────────────────────────────────────────

\subsection{Objectifs}

\begin{table}[H]
\centering
\caption{Périmètre du sprint 7}
\label{tab:s7}
\begin{tabularx}{\textwidth}{|l|L|}
\hline
\rowcolor{headerblue!25}\textbf{Tâche} & \textbf{Résultat attendu} \\\hline
\rowcolor{rowgray}
JasperReports PDF & Rapports congés, formations et absences exportables en PDF
  avec en-tête Synapse et filtre par département/période \\\hline
JasperReports Excel & Export XLSX des mêmes données pour usage tableur \\\hline
\rowcolor{rowgray}
Cache Caffeine & TTL 10 min sur les KPIs analytiques~; invalidation sur
  \texttt{NotificationEvent} Kafka ou appel \texttt{POST /api/admin/attrition/refresh} \\\hline
Module d'attrition & Score de risque (0--100) calculé en Oracle pour chaque
  employé actif~; niveaux FAIBLE/MOYEN/ÉLEVÉ \\\hline
\rowcolor{rowgray}
Tableaux de bord & KPIs multi-rôles~: taux d'absentéisme, répartition des
  demandes, top absents, évolution mensuelle \\\hline
\end{tabularx}
\end{table}

\subsection{Conception~: pipeline analytique}

\begin{figure}[H]
\centering
\begin{tikzpicture}

  \node[infra]     (ora) at (0,3)   {Oracle 23ai\\(requêtes analytiques)};
  \node[component] (rep) at (4.5,5) {ReportService\\JasperFillManager};
  \node[component] (att) at (4.5,1) {AttritionService\\scoring SQL};
  \node[component] (sts) at (9,3)   {StatsService\\cache Caffeine};
  \node[infra]     (kaf) at (9,0.5) {Kafka\\(invalidation)};
  \node[component] (api) at (13,3)  {REST API\\/api/analytics\\:8083};

  \draw[compArr] (ora)  -- (rep);
  \draw[compArr] (ora)  -- (att);
  \draw[compArr] (ora)  -- (sts);
  \draw[compArr] (rep)  -- (api) node[midway,above,font=\scriptsize]{PDF/XLSX};
  \draw[compArr] (att)  -- (api) node[midway,below,font=\scriptsize]{AttritionDTO};
  \draw[compArr] (sts)  -- (api) node[midway,above,font=\scriptsize]{KPIs};
  \draw[compArr,dashed] (kaf) to[bend left=15]
    node[midway,right,font=\scriptsize]{evict cache} (sts);

\end{tikzpicture}
\caption{Pipeline analytique de l'\textit{analytics-service}~: sources, caches et API}
\label{fig:analytics}
\end{figure}

\subsection{Module d'attrition~: modélisation}

\subsubsection*{Facteurs et barème}

\begin{table}[H]
\centering
\caption{Facteurs du score d'attrition et seuils Oracle}
\label{tab:attrition}
\begin{tabularx}{\textwidth}{|l|L|C|}
\hline
\rowcolor{headerblue!25}\textbf{Facteur} & \textbf{Seuil de risque élevé (requête Oracle)} & \textbf{Points max} \\\hline
\rowcolor{rowgray}
Ancienneté & \texttt{MONTHS\_BETWEEN(SYSDATE, HIRE\_DATE) < 6} & 25 \\\hline
Volume d'absences & $\geq 30$ jours de congés approuvés sur les 12 derniers mois & 25 \\\hline
\rowcolor{rowgray}
Taux de rejet & $\geq 50\,\%$ des demandes avec statut \texttt{REJETEE} sur 12 mois & 25 \\\hline
Engagement projets & Aucun enregistrement dans \texttt{PROJECT\_MEMBERS} actif & 25 \\\hline
\end{tabularx}
\end{table}

\subsubsection*{Formule et niveaux de risque}

\begin{equation}
  S_{\text{attrition}} =
  \text{pts}_{\text{ancienneté}} + \text{pts}_{\text{absences}}
  + \text{pts}_{\text{rejets}} + \text{pts}_{\text{engagement}}
  \label{eq:attrition}
\end{equation}

\begin{equation}
  \text{Niveau} =
  \begin{cases}
    \text{ÉLEVÉ}  & \text{si } S \geq 70 \\
    \text{MOYEN}  & \text{si } 40 \leq S < 70 \\
    \text{FAIBLE} & \text{si } S < 40
  \end{cases}
\end{equation}

\subsubsection*{Implémentation~: requête Oracle agrégée}

L'ensemble des quatre facteurs est calculé en une seule passe par une requête
\texttt{WITH} à sous-requêtes \texttt{LEFT JOIN} agrégées, exécutée par
\texttt{AttritionRepository}~. Les résultats sont mis en cache par Caffeine
(TTL 10 min) et invalidés sur appel de
\texttt{POST /api/admin/attrition/refresh} ou sur réception d'un
\texttt{NotificationEvent} de type \texttt{VALIDEE\_RH} dans le flux Kafka.

\begin{lstlisting}[language=SQL,caption={Extrait de la requête d'attrition Oracle}]
WITH anciennete AS (
  SELECT EMPLOYEE_ID,
    CASE WHEN MONTHS_BETWEEN(SYSDATE, HIRE_DATE) < 6 THEN 25 ELSE 0 END pts
  FROM EMPLOYEES WHERE STATUS = 'ACTIF'
),
absences AS (
  SELECT EMPLOYEE_ID,
    CASE WHEN SUM(DAYS_COUNT) >= 30 THEN 25 ELSE 0 END pts
  FROM LEAVE_REQUESTS
  WHERE STATUS = 'APPROUVE'
    AND CREATED_AT > ADD_MONTHS(SYSDATE, -12)
  GROUP BY EMPLOYEE_ID
),
...
SELECT e.EMPLOYEE_ID,
  NVL(a.pts,0) + NVL(ab.pts,0) + NVL(r.pts,0) + NVL(pr.pts,0) AS SCORE
FROM EMPLOYEES e
LEFT JOIN anciennete a  ON a.EMPLOYEE_ID = e.EMPLOYEE_ID
...
\end{lstlisting}

\subsection{Génération de rapports JasperReports}

Les templates \texttt{.jrxml} sont chargés au démarrage depuis
\texttt{classpath:reports/}. La méthode \texttt{ReportService.generate()} reçoit
le type, la période et l'identifiant département, remplit le rapport via
\texttt{JasperFillManager.fillReport(template, params, connection)} et exporte~:

\begin{lstlisting}[language=Java,caption={Export dual PDF/Excel dans ReportService}]
if ("PDF".equals(format)) {
    JasperExportManager.exportReportToPdfStream(print, outputStream);
} else {
    JRXlsxExporter exporter = new JRXlsxExporter();
    exporter.setExporterInput(new SimpleExporterInput(print));
    exporter.setExporterOutput(new SimpleOutputStreamExporterOutput(outputStream));
    exporter.exportReport();
}
\end{lstlisting}

Les rapports sont filtrés par rôle~: un chef accède uniquement aux données de
son département~; l'administrateur RH accède à l'ensemble des données.

\subsection{Résultats du sprint 7}

\begin{figure}[H]
  \centering
  % \includegraphics[width=\textwidth]{images/sprint7_attrition.png}
  \caption{Vue attrition~: KPIs globaux et liste des employés triés par score de risque}
  \label{fig:s7_attrition}
\end{figure}

\begin{figure}[H]
  \centering
  % \includegraphics[width=0.80\textwidth]{images/sprint7_rapport_pdf.png}
  \caption{Rapport de congés généré au format PDF via JasperReports}
  \label{fig:s7_pdf}
\end{figure}

\subsection{Bilan des livrables}

\begin{table}[H]
\centering
\caption{Fonctionnalités livrées à l'issue des sept sprints}
\label{tab:bilan}
\begin{tabularx}{\textwidth}{|L|C|}
\hline
\rowcolor{headerblue!25}\textbf{Fonctionnalité} & \textbf{État} \\\hline
\rowcolor{rowgray}
Infrastructure~: Keycloak, Oracle 23ai, Flyway, Docker Compose & \faCheckCircle \\\hline
Gestion des employés + sync bidirectionnelle Oracle--Keycloak & \faCheckCircle \\\hline
\rowcolor{rowgray}
Workflows demandes RH (5 types, circuit Chef → RH) & \faCheckCircle \\\hline
Workflow prêt avec commission FAVORABLE/DÉFAVORABLE & \faCheckCircle \\\hline
\rowcolor{rowgray}
Notifications asynchrones Kafka + e-mail MailHog & \faCheckCircle \\\hline
Messagerie temps réel WebSocket STOMP/SockJS & \faCheckCircle \\\hline
\rowcolor{rowgray}
Assistant IA conversationnel Groq (\textit{llama-3.3-70b-versatile}) & \faCheckCircle \\\hline
Gestion projets, tâches, Kanban & \faCheckCircle \\\hline
\rowcolor{rowgray}
Recrutement + scoring automatique de CV & \faCheckCircle \\\hline
Rapports JasperReports PDF et Excel & \faCheckCircle \\\hline
\rowcolor{rowgray}
Tableaux de bord analytiques (cache Caffeine, invalidation Kafka) & \faCheckCircle \\\hline
Module de prédiction d'attrition (4 facteurs Oracle) & \faCheckCircle \\\hline
\end{tabularx}
\end{table}

% ─────────────────────────────────────────────────────────────
\section{Conclusion}
% ─────────────────────────────────────────────────────────────

La troisième release complète le périmètre fonctionnel de Synapse~: la gestion
des projets et tâches fournit une mesure objective de l'avancement~; le scoring
Groq automatise la présélection des candidats~; JasperReports produit des documents
exportables adaptés au rôle~; et le module d'attrition calcule directement en Oracle
des indicateurs de risque de départ que les SIRH commerciaux n'auraient pas pu
produire sans développements spécifiques coûteux. Ensemble, ces sept sprints
démontrent qu'une couverture GRH complète est atteignable sur base
exclusivement \textit{open-source} en quatre mois.
